\subsection{優勢}
\begin{enumerate}
    \item \textbf{精細區域劃分}：根據氹仔的地理與功能特徵，將區域分為居民區、旅遊區和樞紐區，針對性地優化站點與投放量，提升了模型的適用性。
    \item \textbf{多因素綜合考慮}：模型整合了空間異質性、時間不對稱性和資源限制等多維約束，使優化結果更貼近實際需求。
    \item \textbf{靈活再平衡策略}：通過引入懲罰成本，允許不完全滿足需求，平衡了運輸成本與用戶體驗，提升了運營靈活性。
    \item \textbf{可擴展性}：支持多輛卡車調度，適應不同規模的運營需求，具有較好的擴展潛力。
\end{enumerate}

\subsection{劣勢}
\begin{enumerate}
    \item \textbf{數據依賴性}：模型效果依賴於歷史需求的準確性，若數據質量不足，優化結果可能失真。
    \item \textbf{計算複雜度}：隨著站點數增加，混合整數規劃與調度模型的求解難度上升，可能需要更高效的算法。
    \item \textbf{參數調節挑戰}：懲罰成本與運輸成本係數等參數需根據實際情況精確設置，否則可能影響結果。
    \item \textbf{動態性不足}：模型以靜態優化為主，對實時需求變化的響應能力有限。
\end{enumerate}
